\scanpage{53}
\mspage{40}
% archive sheet number: 52

objets $C$-cofibrants satisfait aux axiomes de Brown, de plus on \deleted{\unsure{voit}} \inserted{\unsure{prouve}} que
\[
 \mathrm{Ho}_{W}(M_C)\xrightarrow{\approx}\mathrm{Ho}_{W}(M)
\]

est une équivalence. Ainsi les résultats de Brown relatifs à $M_C$, \unsure{muni|munie} des \unsure{ens} de flèches induits par $W,C$, sont applicables à $\mathrm{Ho}_W(M)$ et à ses avatars (tels $\mathrm{Ho}_W({}_S\backslash M)$ ; $\mathrm{Ho}_{\unsure{W|\infty}}(\deleted{\unsure{{}_S\backslash}}M/_{T})$ et $\mathrm{Ho}_{W}({}_S\backslash M_{u_0}/_{T})$ où $u_0:S\unsure{\mathbin{-}|\longrightarrow} T$ est donné.

Il se trouve que pour l'essentiel, les résultats homotopiques \unsure{standart|standarts|standards} (notamment suite \unsure{coexte|coexacte} de cofibration) relatifs aux opérations \unsure{homotopiques} ``à droite'' (ou faut-il dire à gauche ? il s'agit \unsure{pourtant|précisément} des questions d'exactitude à droite, non à gauche) sont valables. Il \deleted{\unsure{me}} faudra revenir dessus de façon plus circonstanciée par la suite. \deleted{Je suis \unsure{particulièrement}} \inserted{Notamment il y} a une bonne théorie des carrés $W$-cocartésiens (et \uline{dans $M_C$}, il y a assez de $W$-cofibrations$^{\times}$, \ldots).

% margin note written sideways
\margin{$^{\times}$plus précisément, par K.S. Brown, \uline{dans} $M_C$ les objets de $C$ sont des $W_C$-cofibrations.}
